\documentclass[11pt]{article}
\usepackage{fontspec}
\usepackage[french]{babel}
\usepackage[leqno]{amsmath}
\usepackage{amssymb}
\usepackage{amsthm}
\usepackage{mathrsfs}
\usepackage{tikz-cd}
\usepackage{enumitem}
\usetikzlibrary{babel}

\newcommand{\printedpage}[1]{\marginpar{\footnotesize p.~#1}}
\newcommand{\msp}[1]{\marginpar{\footnotesize [#1]}}
\newcommand{\M}{\mathcal{M}}
\newcommand{\Fl}{\mathrm{Fl}}
\newcommand{\uFl}{\underline{\mathrm{Fl}}}
\newcommand{\uNerf}{\underline{\mathrm{Nerf}}}

\begin{document}

% ===================== PDF page 1 =====================
\printedpage{---}
% no page number printed on this page
% head: (none)

\begin{center}
{\Large\textsf{LES DÉRIVATEURS}}

\bigskip

\textbf{ALEXANDRE GROTHENDIECK}

\vspace{2cm}

\textbf{Chapitre XV}

\medskip

\textbf{Théorèmes de factorisation}

\textbf{(Modèles (2))}
\end{center}

\clearpage

% ===================== PDF page 2 =====================
\printedpage{---}
% no page number printed on this page
% head: (none)

Ce texte a été déchiffré et transcrit en \LaTeXe{} par M. Künzer. Il a été édité par M. Künzer, J. Malgoire, et G. Maltsiniotis. La transcription est aussi fidèle que possible au manuscrit. Seules quelques corrections mineures évidentes ont été effectuées. Pour les rares ajouts ou commentaires des éditeurs, ainsi que pour la numérotation originale des pages du manuscrit, les caractères de machine à écrire \texttt{[typewriter]} entre crochets sont utilisés.

Cette édition est provisoire. Les remarques, commentaires et corrections sont bienvenus. Envoyer un message à :

\textbf{maltsin@math.jussieu.fr}

\hfill G. Maltsiniotis\hspace*{2cm}

\clearpage

% ===================== PDF page 3 =====================
\printedpage{1}
% head: A. Grothendieck : Dérivateurs, Ch. XV, Édition des Universités de Montpellier II et Paris VII
% foot: Version 5 novembre 2010

\msp{page 1}

\section*{1\quad Construction fondamentale : la suite transfinie de factorisation d'une flèche $\boldsymbol{f}$.}

Soit $\M$ une $\mathfrak{U}$-catégorie, stable par “petites” (\textit{i.e.} $\mathfrak{U}$-petites) $\varinjlim$ filtrantes. Soit
\[
C_0 \;\subset\; \Fl(\M)
\tag*{(1.1)}
\]
un \textit{petit} ensemble de flèches, on se propose de construire pour toute flèche $f$ de $\M$, une factorisation
\[
f \;=\; p\,i \,, \qquad i \;\in\; \widetilde{C_0} = ((C_0)_*)^* \;, \;\; p \;\in\; (C_0)_* \;.
\tag*{(1.2)}
\]
Pour ceci, on va construire une suite ordinale de foncteurs
\[
\left\{
\begin{array}{rcl}
f & \longmapsto & X_\alpha(f) \;=\; X_\alpha(C_0,f) \\
X_\alpha : \uFl(\M) & \longrightarrow & \M
\end{array}
\right.
\tag*{(1.3)}
\]
et des homomorphismes fonctoriels
\[
i_\alpha(f) : \mathrm{source}\, f \;\longrightarrow\; X_\alpha(f) \;,
\tag*{(1.4)}
\]
\[
p_\alpha(f) : X_\alpha(f) \;\longrightarrow\; \mathrm{but}\, f \;,
\tag*{(1.5)}
\]
de sorte qu'on aura le diagramme (où $X = \mathrm{source}\, f$, $Y = \mathrm{but}\, f$)
\[
\begin{tikzcd}[column sep=small]
X \arrow[rr, "f"] \arrow[dr, "i_\alpha(f)"'] & & Y \\
& X_\alpha(f) \arrow[ur, "p_\alpha(f)"'] &
\end{tikzcd}
\qquad ,
\tag*{(1.6)}
\]
ce diagramme étant commutatif
\[
f \;=\; p_\alpha(f) i_\alpha(f) \;.
\tag*{(1.7)}
\]
On peut dire aussi qu'on a une section du foncteur canonique de “composition”, transposé de $\boldsymbol{\Delta}_1 \longrightarrow \boldsymbol{\Delta}_2$, donné par $0 \longmapsto 0$, $1 \longmapsto 2$ :
\[
\left\{
\begin{array}{l}
\begin{tikzcd}[ampersand replacement=\&]
\uNerf_2(\M) \arrow[r, shift left, "\gamma"] \& \uNerf_1(\M) \arrow[l, dotted, shift left, "\Phi_\alpha"]
\end{tikzcd}
\;=\; \uFl(\M) \;, \\[1ex]
\\[1ex]
\\
\end{array}
\right.
\tag*{(1.8)}
\]
$\Phi$ comme “factorisation”,
\printedpage{2}
% head: A. Grothendieck : Dérivateurs, Ch. XV, Édition des Universités de Montpellier II et Paris VII
% foot: Version 5 novembre 2010
\msp{page 2}
où
\[
\Phi_\alpha(f) \;=\; \{\; X \xrightarrow{i_\alpha(f)} X_\alpha(f) \xrightarrow{p_\alpha(f)} Y \;\}
\tag*{(1.9)}
\]
avec $p_\alpha(f)i_\alpha(f) = f$. La fonctorialité de 1.4 et 1.5, ou celle de $\Phi_\alpha$, \textit{i.e.} de 1.9, s'exprime par le fait que tout carré commutatif
\[
\begin{tikzcd}[row sep=large]
X \arrow[r, "f"] \arrow[d, "u_X"'] & Y \arrow[d, "u_Y"] \\
X' \arrow[r, "{f'}"] & Y'
\end{tikzcd}
\tag*{(1.10)}
\]
dans $\M$, interprété comme flèche $u = (u_X, u_Y)$ de $f$ dans $f'$ dans $\uFl(\M)$, donne naissance au double carré commutatif
\[
\begin{tikzcd}[row sep=large]
X \arrow[r, "i_\alpha(f)"] \arrow[d, "u_X"] & X_\alpha(f) \arrow[r, "p_\alpha(f)"] \arrow[d, "X_\alpha(u)"] & Y \arrow[d, "u_Y"] \\
X' \arrow[r, "{i_\alpha(f')}"] & X_\alpha(f') \arrow[r, "{p_\alpha(f')}"] & Y' \;.
\end{tikzcd}
\tag*{(1.11)}
\]
\textbf{N.B.} Quand on voudra expliciter la dépendance de ces constructions par rapport à $C_0$, on écrira
\[
X_\alpha(C_0,f) \;, \qquad i_\alpha(C_0,f) \;, \qquad p_\alpha(C_0,f) \;, \qquad \Phi_\alpha(C_0,f) \;,
\]
ou
\[
X_\alpha(C_0,-) \;\;\text{ou}\;\; X_\alpha^{C_0} \;, \qquad i_\alpha(C_0,-) \;\;\text{ou}\;\; i_\alpha^{C_0} \qquad \textit{etc.} \;,
\]
au lieu de $X_\alpha(f)$, $i_\alpha(f)$, \dots

Voici les propriétés essentielles du système inductif ordinal des $X_\alpha(f)$ ou $\Phi_\alpha(f)$.

\msp{page 3}

\textbf{1.} $X_0(f) = X$, $i_0(f) = \mathrm{id}_X$, $p_0(f) = f$ (\textit{i.e.} la factorisation $\Phi_0(f)$ de $f$ est la factorisation triviale $X \xrightarrow{\mathrm{id}_X} X \xrightarrow{f} Y$).

\textbf{2.}
\[
\begin{aligned}
X_{\alpha+1}(f) &\;=\; X_1(p_\alpha(f)) \\
i_{\alpha+1}(f) &\;=\; i_1(p_\alpha(f)) \circ i_\alpha(f) \\
p_{\alpha+1}(f) &\;=\; p_1(p_\alpha(f))
\end{aligned}
\]
\[
\begin{tikzcd}
\underbrace{X}_{= \,\mathrm{source}\, f} \arrow[r, "i_\alpha(f)"] \arrow[rrr, bend left=20, "f"] & X_\alpha(f) \arrow[rr, "p_\alpha(f)"] \arrow[dr, bend right=40, "i_1(p_\alpha(f))"'] & & Y \\
& & \underbrace{X_1(p_\alpha(f))}_{= X_{\alpha+1}(f)} \arrow[ur, bend right=40] &
\end{tikzcd}
\]

\clearpage

% ===================== PDF page 5 =====================
\printedpage{3}
% head: A. Grothendieck : Dérivateurs, Ch. XV, Édition des Universités de Montpellier II et Paris VII
% foot: Version 5 novembre 2010

ou mieux
\[
\begin{tikzcd}[row sep=large]
X \arrow[rrr, "f"] \arrow[d, "\mathrm{id}"] & & & Y \arrow[d, "\mathrm{id}"] \\
X \arrow[r, "i_\alpha(f)"] \arrow[d, "\mathrm{id}"] & X_\alpha(f) \arrow[rr, "p_\alpha(f)"] \arrow[d, "\mathrm{id}"] & & Y \arrow[d, "\mathrm{id}"] \\
X \arrow[r] & X_\alpha(f) \arrow[r, "i_1(p_\alpha(f))"] & \underbrace{X_1(p_\alpha(f))}_{= X_{\alpha+1}(f)} \arrow[r, "p_1(p_\alpha(f))"] & Y \;.
\end{tikzcd}
\tag*{(1.12)}
\]
Le morphisme de transition $X_\alpha(f) \longrightarrow X_{\alpha+1}(f)$ est
\[
X_\alpha(f) \;\; \xrightarrow{\;\;i_1(p_\alpha(f))\;\;} \;\; X_1(p_\alpha(f)) \;=\; X_{\alpha+1}(f) \;.
\]

\textbf{3.} Si $\alpha$ est un ordinal limite, on a
\[
\left\{
\begin{array}{rcl}
X_\alpha(f) & = & \varinjlim\limits_{\alpha'<\alpha} X_{\alpha'}(f). \\[3ex]
i_\alpha(f) & = & \parbox[t]{8cm}{morphisme canonique de $X = X_0(f)$ dans la limite inductive.} \\[5ex]
p_\alpha(f) & = & \text{limite projective des } p_{\alpha'}(f) : X_{\alpha'}(f) \longrightarrow Y.
\end{array}
\right.
\]
(On peut considérer la limite inductive définissant $X_\alpha(f)$ comme une limite dans $X\backslash\M/Y\dots$)

Les morphismes de transition
\[
X_{\alpha'}(f) \;\longrightarrow\; X_\alpha(f) \qquad (\; \alpha' \;<\; \alpha \;)
\]
sont les morphismes canoniques relatifs au système inductif des $X_{\alpha'}(f)$.

\msp{page 4}

On peut donc dire que la suite ordinale des factorisations fonctorielles $\Phi_\alpha(f)$ de $f$ variable, se déduit de la factorisation fonctorielle $\Phi_1(f)$ de $f$ en
\[
\begin{tikzcd}
X \arrow[r, "i_1(f)"] \arrow[rr, bend left=20, "f"] & X_1(f) \arrow[r, "p_1(f)"] & Y
\end{tikzcd}
\]
par “itération transfinie à droite”, \textit{i.e.} en regardant les opérations $\Phi_\alpha$ comme des foncteurs
% footnote mark (1) printed beside the following display; \footnotemark/\footnotetext used because \footnote does not work inside math
\[
\left\{
\begin{array}{rcl}
p_\alpha :\; \uFl(\M) & \longrightarrow & \uFl(\M) \\
f & \longmapsto & p_\alpha(f) \;,
\end{array}
\right.
\qquad\qquad \text{\footnotemark}
\]
\footnotetext{\textbf{N.B.} $p_\alpha$ respecte les structures de catégories \textit{sur} $\M$, via les foncteurs but : $\uFl(\M) \longrightarrow \M$, \textit{i.e.} $\uFl(\M)$ est considérée comme objet de $\mathrm{CAT}/\M$.}
munis d'homomorphismes fonctoriels (au-dessus de $\M$)
\[
\mathrm{id}_{\uFl(\M)} \;\longrightarrow\; p_\alpha \;.
\]

\end{document}
